\documentclass[11pt]{amsart}
\usepackage{amsmath,amssymb,amsthm,mathrsfs}
\usepackage[margin=1in]{geometry}
\usepackage{enumitem}
\usepackage{hyperref}

\newcommand{\cA}{\mathcal{A}}
\newcommand{\cM}{\mathcal{M}}
\newcommand{\cW}{\mathcal{W}}
\newcommand{\barB}{\overline{B}}
\newcommand{\Ran}{\mathrm{Ran}}
\newcommand{\MC}{\mathrm{MC}}
\newcommand{\Hom}{\mathrm{Hom}}
\newcommand{\Ext}{\mathrm{Ext}}
\newcommand{\Conf}{\mathrm{Conf}}
\newcommand{\FM}{\mathrm{FM}}
\newcommand{\Tr}{\mathrm{Tr}}
\newcommand{\Sym}{\mathrm{Sym}}
\newcommand{\End}{\mathrm{End}}
\newcommand{\g}{\mathfrak{g}}
\newcommand{\h}{\mathfrak{h}}
\newcommand{\n}{\mathfrak{n}}
\newcommand{\fsl}{\mathfrak{sl}}
\newcommand{\Z}{\mathbb{Z}}
\newcommand{\Q}{\mathbb{Q}}
\newcommand{\R}{\mathbb{R}}
\newcommand{\C}{\mathbb{C}}

\theoremstyle{plain}
\newtheorem{theorem}{Theorem}[section]
\newtheorem{proposition}[theorem]{Proposition}
\newtheorem{conjecture}[theorem]{Conjecture}
\newtheorem{problem}[theorem]{Problem}
\theoremstyle{definition}
\newtheorem{definition}[theorem]{Definition}
\theoremstyle{remark}
\newtheorem{remark}[theorem]{Remark}

\title{Frontier Research Questions\\[4pt]
\large Modular Koszul Duality for Chiral Algebras on Curves}
\author{Raeez Lorgat}
\date{April 2026}

\begin{document}
\maketitle

\tableofcontents

\section{Introduction}

The programme of modular Koszul duality studies the bar-cobar
adjunction for chiral algebras on algebraic curves, with the
Maurer--Cartan element
$\Theta_\cA \in \MC(\g_\cA^{\mathrm{mod}})$ as the universal
organising structure. The theorem packages A, B, C, D, and H
have proved cores, with the stated qualifiers on the C2 and
Hochschild boundary cases.  The shadow obstruction tower is
constructed on the verified lanes, and the genus expansion
$F_g = \kappa(\cA) \cdot \lambda_g^{\mathrm{FP}}$ is established
on the uniform-weight lane.  For multi-weight algebras one must
add the cross-channel correction
$\delta F_g^{\mathrm{cross}}(\cA)$, with
$\delta F_1^{\mathrm{cross}}=0$.  This document records the
frontier questions that remain, with mathematical context,
physical motivation, and available tools.

The standard is that of Beilinson: a claim is false until
independently verified; a coincidence is an error until explained;
a narrative is dangerous until every sentence is deducible from
the mathematics.

\section{The coderived BV/BRST identification}

\subsection{The problem}

Let $\cA$ be a chirally Koszul algebra on a smooth curve $X$,
and let $\Sigma_g$ denote a closed Riemann surface of genus~$g$.
The Batalin--Vilkovisky complex $\mathrm{Obs}(\Sigma_g, \cA)$ of
the two-dimensional sigma model with target the chiral algebra
$\cA$ and the bar complex $B^{(g)}(\cA)$ of the genus-$g$
factorisation homology are expected, after choosing the correct
realisation and pairing, to compute the same observables.  The
theorem here is the algebraic comparison of the complexes, not the
full amplitude extraction.

At genus~$0$, the identification $\mathrm{Obs}(\Sigma_0, \cA)
\simeq B^{(0)}(\cA)$ is consistent with the Costello--Gwilliam
BV/factorisation formalism and is proved here as the genus-$0$
chiral-bar comparison. At
genus $g \geq 1$, the bar complex acquires curvature
$m_0 = \kappa(\cA) \cdot \omega_g$ from the Kodaira--Spencer
deformation, and the BV complex acquires the BV Laplacian
$\Delta_{\mathrm{BV}}$ as a genus-raising operator.

\subsection{The resolution}

The quartic harmonic discrepancy between the two complexes has
the universal form
\[
\delta_4^{\mathrm{harm}}
= c_4(\cA)\,m_0
\]
when $\kappa\neq0$ and the harmonic propagator is compared with
the curvature $m_0=\kappa\omega_g$.  For Virasoro,
$c_4(\mathrm{Vir})=10/(c(5c+22))$.  In the coderived category
$D^{\mathrm{co}}(\cA)$ of Positselski, the curvature-divisible
comparison cone is coacyclic.  Thus $\delta_4$ is not an ordinary
coboundary of the holomorphic bar complex; the comparison becomes
an isomorphism after passage to $D^{\mathrm{co}}$. Higher-arity discrepancies
$\delta_r^{\mathrm{harm}} \propto c_r(\cA)m_0^{\lfloor r/2\rfloor-1}$
are absorbed by the same curvature filtration.

\begin{theorem}[BV/bar in $D^{\mathrm{co}}$]
For every chirally Koszul algebra~$\cA$, the BV and bar
complexes are identified in $D^{\mathrm{co}}(\cA)$ at every
genus.
\end{theorem}

At the chain level, the identification holds unconditionally for
shadow-depth classes $\mathsf{G}$ and $\mathsf{L}$, holds for
class~$\mathsf{C}$ under harmonic-decoupling hypotheses, and
fails in the ordinary holomorphic chain category for
class~$\mathsf{M}$ (Virasoro, $\cW_N$): the
harmonic obstruction $1/\mathrm{Im}(\tau)$ is non-holomorphic
and hence not a coboundary of the holomorphic bar differential.
The coderived category is a canonical categorical resolution in
which the identification is universal.  Equivalent pro-object,
$J$-adic, and filtered weight-completed chain ambients also exist;
the bounded direct-sum chain category is the ambient that fails.

\subsection{The three viewpoints}

The inevitability of the coderived category can be understood from
three complementary perspectives.

\subsubsection{The worldsheet perspective}

On a genus-$g$ surface~$\Sigma_g$, the sewing construction
introduces the propagator $P = P_{\mathrm{bar}} +
P_{\mathrm{exact}} + P_{\mathrm{harm}}$. The bar complex uses
$P_{\mathrm{bar}} = d\log E(z,w)$ (meromorphic, weight~$1$);
the BV complex uses the full~$P$. The discrepancy at arity~$r$
is $\delta_r = c_r(\cA) \cdot P_{\mathrm{harm}}^{\lfloor
r/2 \rfloor - 1}$. The identity
$P_{\mathrm{harm}} = m_0/\kappa$ gives
$\delta_4 = (Q^{\mathrm{contact}}/\kappa) \cdot m_0$.
The discrepancy is proportional to the curvature. In
$D^{\mathrm{co}}$, where coacyclic curvature cones are inverted,
these modes vanish in the coderived comparison.  Physical amplitude
extraction remains conditional on the chosen realisation and
pairing.

\subsubsection{The holographic perspective}

In the Costello--Li framework, the bulk theory on
$\mathrm{AdS}_{d+1}$ uses Witten diagrams with the
bulk-to-boundary propagator; the boundary theory uses the bar
complex with the OPE. At one loop, the mismatch is the
holographic anomaly: the Weyl anomaly coefficient becomes
$\delta_4 = Q^{\mathrm{contact}} \cdot m_0$ in the algebraic
framework. The coderived resolution says: this anomaly is
\emph{exact} in the correct cohomological sense. Holography is
exact at the level of cohomology, not representatives.

\subsubsection{The categorical perspective}

$D^{\mathrm{co}}(\cA)$ is the homotopy category of dg
comodules over $B(\cA)$ with $d^2 = m_0 \cdot \mathrm{id}$
permitted. The bar complex $B(\cA)$ satisfies $D_B^2 = 0$
always (curvature absorbed into the total differential). The BV
complex satisfies $d^2_{\mathrm{BV}} = \kappa \cdot \omega_g$,
the same curved relation. The coderived quasi-isomorphism
identifies the two in the \emph{minimal} category where both are
well-defined, without auxiliary strictification choices.

\subsection{The class-by-class mechanism}

The chain-level analysis reveals a hierarchy:
class~$\mathsf{G}$ (no vertices: $\delta_r = 0$),
$\mathsf{L}$ (Jacobi kills harmonic coupling),
$\mathsf{C}$ (role separation: fundamental generators $\neq$
quartic source), $\mathsf{M}$ (triple role of~$T$ prevents
decoupling). For $\mathsf{M}$, the stress tensor simultaneously
generates the bar complex, sources $Q^{\mathrm{contact}}$, and
is contracted by~$\Delta_{\mathrm{BV}}$: no algebraic identity
can kill the coupling. The coderived resolution works at a
different level: not by killing $\delta_4$ but by recognising it
as $Q^{\mathrm{contact}} \cdot m_0$, a curvature multiple that
is exact in~$D^{\mathrm{co}}$.

\subsection{Physical significance}

The coderived resolution proves the algebraic BV/bar comparison:
BV quantisation and bar construction produce quasi-isomorphic
complexes in $D^{\mathrm{co}}$.  The tree-level amplitude pairing
and full amplitude extraction remain conditional on the physical
realisation.  The curvature absorption is the algebraic model for
the principle that higher-genus UV divergences should be absorbed
by the correct categorical ambient. The
theorem also explains why the scalar genus expansion
$F_g = \kappa \cdot \lambda_g^{\mathrm{FP}}$ on the
uniform-weight lane agrees with the BV effective action:
both are representatives of the same $D^{\mathrm{co}}$-cohomology
class.  Multi-weight algebras require the additional
cross-channel term $\delta F_g^{\mathrm{cross}}$.


\section{Non-principal DS-KD commutation: the (3,2) partition}

\subsection{Context}

Drinfeld--Sokolov reduction $\mathrm{DS}_f\colon V_k(\g) \to
W_k(\g, f)$ for nilpotent $f \in \g$ produces $W$-algebras from
affine Kac--Moody algebras. Koszul duality
$\cA \mapsto \cA^!$ produces the dual algebra. The fundamental
question: does DS commute with Koszul duality?

\subsection{What is proved}

\begin{enumerate}[label=(\roman*)]
\item Principal nilpotent $f_{\mathrm{princ}}$ (all types): proved.
\item Hook-type partitions $[N{-}1, 1]$ in type~$A$:
 conditional transport, with the compatibility criteria proved
 and the full DS/bar compatibility hypothesis still explicit
 (Fehily--Creutzig--Linshaw--Nakatsuka--Sato, transport
 propagation along the Hasse diagram).
\item Self-transpose rectangular $(2,2)$ in $\fsl_4$: proved by
 self-duality ($c + c' = 14$, self-dual at $c^* = 7$).
\end{enumerate}

\subsection{The open case}

\begin{problem}
Let $f = f_{(3,2)}$ in $\fsl_5$. The partition $(3,2)$ is
non-self-transpose (its transpose is $(2,2,1)$) and non-hook.
The BRST complex has non-abelian $\n_+$ and the first genuinely
non-hook ghost-sector pattern. Does
$\mathrm{DS}_{f_{(3,2)}}(V_k(\fsl_5)^!) \simeq
\mathrm{DS}_{f_{(3,2)}}(V_k(\fsl_5))^!$?
\end{problem}

The correct nilradical datum is
$\dim\n_+=10$, with half-integer grades $1/2,1,3/2,2$,
four-step nilpotence, ten nonzero bracket families, and ghost
sectors $-10,\ldots,10$.  The number~$8$ is the
centralizer/strong-generator count, not $\dim\n_+$.  The
computation requires sparse BRST matrices decomposed by the
$21$ ghost sectors.  This is the first non-hook non-abelian test
of DS-KD commutation.

\subsection{Physical motivation}

Non-principal $W$-algebras arise as boundary algebras of 3d
$\mathcal{N}=4$ gauge theories with non-maximal Nahm pole
boundary conditions (Gaiotto--Witten). A proof of DS-KD
commutation would establish that Koszul duality of the boundary algebra is
compatible with varying the Nahm pole type: a form of ``duality
of boundary conditions'' in holography.


\section{The genus-5 cross-channel correction}

\subsection{The candidate universal $N$-formula}

For the principal $W$-algebra $\cW_N$ with generators at
conformal weights $2, 3, \ldots, N$, the genus-$g$ free energy
decomposes as
\[
F_g(\cW_N, c)
= \kappa(\cW_N) \cdot \lambda_g^{\mathrm{FP}}
+ \delta F_g^{\mathrm{cross}}(\cW_N, c),
\]
where the cross-channel correction $\delta F_g^{\mathrm{cross}}$
is computed in low genus and may admit universal polynomial
skeleta in~$N$, vanishing at $N = 2$ (uniform weight). At
genus~$2$ the gravitational skeleton is
$\delta F_2^{\mathrm{grav}} = B_2(N) + A_2(N)/c$ with
$A_2(N) = (N{-}2)(3N^3{+}14N^2{+}22N{+}33)/24$.
This is exact for $N=3$ and a lower bound for $N\geq4$.
At genus~$3$: four terms with verified polynomial degrees
$(1,3,5,7)$ in~$N$.  At genus~$4$, $d_4=4$ is the proved
$\cW_3$ numerator degree in~$c$, not a universal $N$-degree
statement.

\begin{problem}
Compute $\delta F_5^{\mathrm{cross}}(\cW_N, c)$ and test whether
a six-term universal $N$-formula exists.  For $\cW_3$, verify the
predicted $c$-polynomial pattern $d_5 = 5$, $n_5 = 6$ terms, net
degree~$1$.
Test whether a well-defined $A_{\mathrm{cross}}/A_{\mathrm{scalar}}$
exists; if single-action Borel control is proved, compute it.
\end{problem}

This requires enumerating all ${\sim}4000$ stable graphs of
$\overline{\cM}_{5,0}$ and computing the gravitational
mixed-channel amplitude for each.

\subsection{Physical significance}

The cross-channel correction quantifies the contribution of
mixed-propagator graphs: worldsheet configurations where
different OPE channels interact. These have no analogue in
the scalar (single-channel) theory and represent genuinely new
multi-channel structure.  If a single cross-channel action exists,
the ratio $A_{\mathrm{cross}}/A_{\mathrm{scalar}}$ would measure
how much ``heavier'' the cross-channel growth is than the scalar
one.  At present this is a conjectural asymptotic diagnostic, not a
universal invariant.


\section{Admissible Koszulness at rank $\geq 2$}

\subsection{The rank obstruction}

\begin{conjecture}[Admissible Koszul rank obstruction]
For every simple Lie algebra $\g$ of rank $r \geq 3$ and every
admissible level $k = p/q - h^\vee$ with denominator $q \geq 3$,
the Li--bar $E_2$ page contains a rank-$r$ Cartan obstruction to
PBW concentration.
\end{conjecture}

The obstruction arises from the abelian Cartan subalgebra
$\h \subset \g$: the Poisson differential $d_1$ on the Li--bar
$E_1$ page kills all root generators
(via $[e_\alpha, f_\alpha] = h_\alpha$) but the Cartan
generators $h_1, \ldots, h_r$ survive because $\h$ is abelian.

At rank~$1$ ($\fsl_2$): the Cartan is $1$-dimensional, but the
PBW spectral sequence has at most $2$ non-zero columns
($\dim \fsl_2 = 3$), forcing all higher differentials to vanish.
Hence $L_k(\fsl_2)$ is Koszul at all admissible levels.

At rank~$2$ ($\fsl_3$): the obstruction is proved for
$q \geq 3$, with classes appearing at $h_\alpha=(p-1)q$.
The $q\leq 2$ cases are compatible with the same picture but
remain outside the proved range.  At rank~$3$ ($\fsl_4$):
the predicted obstruction has rank~$3$ and the computation is
feasible.

\subsection{Physical significance}

Admissible-level simple quotients are the VOAs of rational
conformal field theories (Arakawa).  For $\fsl_3$, bar-cobar
inversion fails at $q\geq3$ by theorem.  For rank $\geq3$, the
same obstruction is conjectural: the null vectors that make the
theory rational are expected to obstruct PBW concentration.  This would be the
algebraic shadow of the representation-theoretic rigidity of
rational CFTs. (At rank~$1$, $L_k(\fsl_2)$ is Koszul at all
admissible levels.)


\section{DK-4 and DK-5: Extending the Drinfeld--Kohno bridge}

\subsection{DK-4: Beyond the evaluation core}

MC3 gives the proved evaluation-core reduction for simple types.
DK-4 asks whether that core assembles into the full global
algebraic identification; thick generation beyond the proved core is
part of the remaining problem.

\begin{problem}[DK-4]
Is the evaluation-generated core of the DK category dense in the
full module category of $U_q(\g)$, in the sense that the
localisation functor is fully faithful on ind-completions?
\end{problem}

The gap is algebraic: pointwise Ext computations (at individual
modules) do not automatically globalise to the category level.
For $\fsl_3$, genuinely non-evaluation modules (Whittaker,
Wakimoto) exist and must be incorporated.

\subsection{DK-5: Categorical $E_1$ primacy}

\begin{problem}[DK-5]
Does the ordered factorisation structure on $\Ran(X)$ induce a
braided monoidal equivalence between the ordered bar category and
the representation category of the Yangian $Y(\g)$, not just on
the evaluation core but on the full factorisation category?
\end{problem}

This is the categorical incarnation of $E_1$ primacy: the
Yangian and its braided monoidal structure are native to the
ordered bar complex, and the passage to the full $\Ran$-space
category requires the Francis--Gaitsgory pro-nilpotent
completion.


\section{Grand completion}

\begin{problem}
Construct a Quillen equivalence (or $\infty$-categorical
equivalence) between the category of chiral algebras on~$X$ and
the category of conilpotent chiral coalgebras on~$X$, with
bar-cobar as the Quillen pair.
\end{problem}

This requires: (a) model category structures on both sides,
(b) a cumulant recognition theorem (characterising the image of
the bar functor among all coalgebras), and (c) a jet principle
(extending formal-disk data to global curve data within the model
structure). This is the hardest problem in the programme.


\section{Analytic realisation: three-layer gap}

Three frameworks address the non-perturbative completion of the
formal genus expansion:
\begin{enumerate}[label=(\alph*)]
\item The sewing envelope $\cA^{\mathrm{sew}}$: Hausdorff
 completion for sewing-amplitude seminorms. Proved in the
 Heisenberg and lattice cases; polynomial OPE growth plus
 subexponential sector growth is the sufficient condition
 expected to cover the interacting standard landscape, but that
 analytic verification remains conditional.
\item Conformally flat factorisation homology in
 $\mathrm{IndHilb}$ (Moriwaki26a): analytic, $1$-categorical,
 metric-dependent; the Heisenberg/Bergman $2$-disk example is the
 separate Moriwaki26b lane.  Metric-independence and
 non-Heisenberg extensions remain conditional.
\item Coderived shadow invariants $Q_g^{\mathrm{an}}(\cA)$:
 categorical, using Positselski's coderived category with
 curvature $d^2 = m_0 \cdot \mathrm{id}$.
\end{enumerate}

\begin{problem}
Unify these three frameworks into a single ``analytic bar
complex'' $B^{\mathrm{an}}(\cA)$ that simultaneously carries the
algebraic bar differential, the analytic sewing data, and the
coderived curvature.
\end{problem}


\section{$E_1$ Verdier duality on ordered configurations}

The Verdier intertwining $\mathbb{D}_{\Ran}(\barB(\cA))
\simeq \barB(\cA^!)$ operates on the symmetric bar over
unordered $\Ran(X)$. The ordered bar $\barB^{\mathrm{ord}}(\cA)$
lives on $\Conf^<_n(X)$, which is \emph{not} a factorisation
space. Therefore $\mathbb{D}_{\Ran}(\barB^{\mathrm{ord}})$ is
undefined.

\begin{problem}
Develop a ``ribbon $\Ran$'' framework in which Verdier duality
extends to ordered configurations, and establish the ordered
Verdier intertwining
$\mathbb{D}_{\mathrm{rib}}(\barB^{\mathrm{ord}}(\cA))
\simeq \barB^{\mathrm{ord}}((\cA^!)^{\mathrm{op}})$.
\end{problem}

The correct $E_1$ analogue uses opposite-coalgebra duality
(reversal involution), not Verdier duality. The ribbon $\Ran$
development would provide the geometric home for the ordered
Koszul dual.


\section{Resurgence: the cross-channel instanton action}

The scalar instanton action $A_{\mathrm{scalar}} = (2\pi)^2$ is
proved in the uniform-weight/scalar lane.  Low-genus cross-channel
data suggest faster growth, but the cross-channel asymptotics are
not yet controlled by a single universal instanton action in the
usual Borel sense.

\begin{problem}
Determine whether a well-defined
$A_{\mathrm{cross}}/A_{\mathrm{scalar}}$ exists and, if so, compute
it.  The current interval $[1.7, 3.1]$ is a three-data-point
extrapolation; genus-$5$ data would test rather than prove
single-action control.
\end{problem}

If a single cross-channel action exists, it should come from a
non-perturbative configuration in which different OPE channels
interact coherently across the worldsheet.  No such saddle has yet
been constructed.


\section{Cross-channel generating function}

\begin{problem}
Does the cross-channel genus expansion
$\sum_{g \geq 2} \delta F_g^{\mathrm{cross}}(\cW_N, c)\,
\hbar^{2g}$ admit a closed-form generating function in the
bivariate $(g, N)$?
\end{problem}

For $\cW_3$, the exact $g=2,3,4$ data rule out the listed
one-variable, $\hat A$-type, geometric, and algebraic ansatz
classes.  The broader bivariate $\cW_N$ closed-form problem remains
open.  Low-genus data
suggest non-separability: the $N$-dependence at genus~$g$ involves
high-degree power sums of generator weights. A plausible framework
is the $\cW_N$ Frobenius manifold with $N{-}1$ flat coordinates,
rather than a one-variable $\hat{A}$-genus series.


\section{Scalar saturation residual}

\begin{conjecture}[Scalar saturation]
For a simple-factor modular Koszul algebra~$\cA$,
$\dim H^2_{\mathrm{cyc}}(\cA) = 1$.  More generally, with
$r$ simple-current factors, the expected dimension is~$r$; the
stronger residual assertion is that the resulting classes satisfy
$\Gamma_i=\kappa_i\Lambda$ on each factor.
\end{conjecture}

A violation would produce ``isomodular'' algebras: chiral
algebras with identical $\kappa$ but different higher-genus
amplitudes. No counterexample is known among $61+$ tested
families across five independent verification paths.

\bigskip
\hrule
\bigskip

\noindent\textit{Status.  Except for explicitly cited proved inputs,
these are open or conjectural extensions beyond the specified proved
lanes.  Each requires a separate proof or computation before it can
be promoted to theorem status.}

\end{document}
